% PDF707 begins here at Natural experiment; the sourcepage-707 label is in A-M.
\glossentry{Natural experiment}{自然实验}{一种机会性研究设计，其中接受某种治疗或暴露于某危险因素很可能是随机的，或者至少与结局的决定因素无关，因此有助于估计因果效应。例如，美洲原住民拥有的赌场开业，被用作自然实验，研究收入增加对多种结局的影响。（总体而言，这些影响是有利的！）}

\glossentry{Negative predictive value}{阴性预测值}{见\emph{预测值，阴性}。}

\glossentry{Nested case-control study}{巢式病例对照研究}{病例与对照从一个已明确界定的较大队列，或先前已经入组的队列研究参与者中选取。对队列所有人进行某些测量过于昂贵时，通常采用这种设计。例如，研究者计划用巢式病例对照研究，确定新生儿筛查血斑中细胞因子的相对水平，是否与后来发生脑瘫有关；测量将覆盖所有脑瘫婴儿的样本，以及从对照中随机抽取的1\%样本。}

\sourcepage{708}
\glossentry{Nested double (or multiple) cohort study}{巢式双（或多）队列研究}{从较大队列中，分别抽取暴露者与未暴露者，或不同程度暴露者组成独立队列。暴露罕见时，巢式多队列研究比随访整个队列更有效率。例如，在一个大型医疗系统中，可以比较接受金属对聚乙烯关节面髋关节置换者，与接受陶瓷对聚乙烯植入物者的结局。}

\glossentry{No-cost extension}{不追加经费延期}{研究资助期结束后，仍可使用预算结余，完成或延续资助决定中所规定工作的时期。例如，美国国立卫生研究院（NIH）批准六个月的不追加经费延期，让研究者完成数据分析和论文撰写。}

\glossentry{N-of-1 design}{单例设计}{只纳入一人的交叉试验，也称单患者试验。患者在规定时期内随机接受盲法治疗、安慰剂或另一有效治疗，随后转为另一种干预，持续同样长的时间。例如，一名临床医师与患者采用单例试验，判断在触发点注射利多卡因，缓解疼痛是否比注射生理盐水更有效。}

\glossentry{Nominal variable}{无序分类变量}{类别之间没有内在顺序的分类变量。例如，把宗教信仰归为佛教、基督教、印度教、犹太教、伊斯兰教、其他或无宗教，并编码为无序分类变量。另见\emph{有序分类变量}。}

\glossentry{Nondifferential misclassification bias}{非差异性错分偏倚}{不受参与者是病例还是对照影响，或有时不受其是否暴露影响的一类偏倚。非差异性错分通常会缩小表面的组间差异，使关联更难发现。例如，病例和对照对既往抗生素暴露的回忆都不完全准确，但回顾病历发现两组不准确程度相近，因此这一偏倚看来是非差异性的。另见\emph{差异性错分偏倚}。}

\glossentry{Noninferiority margin}{非劣效界值}{预先规定的两种治疗疗效差，常称delta（$\Delta$）。这一差异足够小，使研究者可以认为新治疗并不明显劣于标准治疗。例如，在新抗凝药与华法林预防房颤患者卒中的非劣效试验中，界值设为卒中率相差1个百分点。}

\glossentry{Noninferiority trial}{非劣效试验}{比较一种相对已有治疗具有某些优点的新治疗，例如更安全、更便宜或更易使用，旨在证明其疗效不劣于已有治疗的临床试验。
\sourcepage{709}
例如，一项研究不引起嗜睡的新镇痛药的非劣效试验，证实其缓解术后疼痛的作用不劣于羟考酮。}

\glossentry{Nonprobability sample}{非概率样本}{无法确定可接触人群中某成员被纳入可能性的样本。许多，甚至大多数临床研究样本都属于此类。例如，研究者只要有人能够访谈，就把骨科诊所就诊患者纳入便利样本，通常在不太忙的日子进行；这就是一种非概率样本。另见\emph{概率抽样}。}

\glossentry{Nonresponse bias}{未应答偏倚}{未作回答，例如未回答问卷，影响研究结果的一类偏倚。例如，研究非法药物使用与肾衰竭风险时，关于用药的问题常被留空，因此研究者担心未应答偏倚。}

\glossentry{Normalization}{规范化}{使关系数据库遵守一组规则的过程，这些规则有助于查询、维护参照完整性，并减少冗余和内部不一致。规范化常把宽表中重复的列，转换为窄表中的独立行。例如，一个未规范化的Participant表，为多次实验室检查设置labID1、labresult1、labdate1、labID2、labresult2、labdate2等字段。规范化时，应建立新的LabResult表，包含主键、labID、labresult、labdate，以及连接回参与者个人资料的外键（见图19.3）。}

\glossentry{Null hypothesis}{零假设}{规定所比较组别间没有差异的一种研究假设形式。例如，零假设规定：在血脂正常的参与者中，他汀组与安慰剂组发生间歇性跛行的风险相同。}

\glossentry{Number needed to harm}{需伤害人数}{为了造成额外一个结局发生，需要接受某种治疗的绝对人数，以风险差的倒数计算。例如，雌激素使用所致静脉血栓栓塞事件的风险差为0.3\%，需伤害人数就是333。另见\emph{需治疗人数}。}

\sourcepage{710}
\glossentry{Number needed to treat}{需治疗人数}{为了预防一个结局发生，需要接受治疗的绝对人数，以绝对风险降低的倒数计算。例如，评价妇女、婴幼儿特殊补充食品计划（WIC）的利益时，大约需要让25名孕妇参加，才能预防一例低出生体重婴儿。}

\glossentry{Numeric variable}{数值变量}{能够用数字量化的变量。例如，用kg/m$^2$表示的体重指数（BMI），就是数值变量。}

\glossentry{Observational study}{观察性研究}{研究者不实施任何干预，只观察参与者的一类研究设计的总称，因此不包括随机试验。例如，研究者开展观察性研究，确定黑色素瘤的危险因素。}

\glossentry{Observer bias}{观察者偏倚}{研究者或研究助理知道参与者的一项或多项特征，例如是病例还是对照，或者是否暴露于特定危险因素，因而作出非客观评价的情况。例如，访谈时，西班牙裔青少年比亚裔更容易被认为存在愤怒管理问题；但自填调查及学校记录回顾未发现两组不同，因此访谈发现似乎由观察者偏倚造成。}

\glossentry{Observer variability}{观察者变异}{实施测量的观察者存在变化，导致同一变量重复测量不同。例如，记录腰围的检查者可能把皮尺放在腹部不同高度，造成测量的观察者变异。}

\glossentry{Odds}{比值}{疾病或其他结局的风险，除以$1-\text{风险}$。例如，高风险女性终生乳腺癌风险为20\%，则终生发病比值为$0.20/0.80=0.25$。罕见疾病，即发生于不到约10\%人群的疾病，其风险与比值相近。}

\glossentry{Odds ratio}{比值比}{暴露于某危险因素者的疾病或其他结局比值，除以未暴露者相应比值；也等于有病者，例如病例对照研究中的病例，其暴露比值除以无病者，即对照的暴露比值。疾病在暴露及未暴露组都罕见时，因风险与比值相近，风险比与比值比也相近。例如，
\sourcepage{711}
美国终末期肾病每年发生于约0.04\%人口；一项病例对照研究中，非胰岛素依赖型糖尿病的比值比为7，意味着这些患者发生终末期肾病的可能性约为无糖尿病者的七倍。（Perneger TV, Brancati FL, Whelton PK, Klag MJ. End-stage renal disease attributable to diabetes mellitus. Ann Intern Med. 1994;121:912-918.）}

\glossentry{One-sample t test}{单样本$t$检验}{比较样本中某变量均值与一个固定常数的统计检验。最常见的单样本$t$检验是配对$t$检验，即将配对测量差值的样本均值与零比较，例如同一参与者在不同时间的测量。例如，研究者发现，男性在住院医师培训期间平均增重$4\pm3$ kg（均值$\pm$标准差），单样本$t$检验$P=0.03$。另见\emph{双样本$t$检验}。}

\glossentry{One-sided hypothesis}{单侧假设}{研究者只关注两个可能方向之一的I类错误可能性的备择假设，例如只关注治疗使结局风险升高，或只关注使风险降低，而非同时关注两个方向。例如，研究者检验了“吸烟与痴呆风险增加有关”这一单侧假设。另见\emph{双侧假设}。}

\glossentry{One-to-many}{一对多}{数据库中“一个”参与者或属性，对应“多个”相关测量的一种关系。例如，一项研究反复测量体重，每名参与者都有初级保健服务提供者及居住城市。参与者可处于“一”侧，体重测量处于“多”侧；城市可处于“一”侧，参与者及保健服务提供者处于“多”侧；保健服务提供者也可处于“一”侧，其患者处于“多”侧。另见\emph{关系数据库}。}

\glossentry{Open-ended question}{开放式问题}{旨在让人以自由形式或开放文字作答的问题。例如，研究者问：“您认为我们社区面临的最大健康威胁是什么？”让参与者用自己的话回答。另见\emph{封闭式问题}。}

\glossentry{Operations manual}{操作手册}{研究方案的扩展版本，通常包括方案本身、研究组织与政策信息，以及详细的方法部分。例如，手册规定，参与者须安静坐五分钟后才测血压。}

\sourcepage{712}
\glossentry{Opportunistic study design}{机会性研究设计}{因某个特殊机会而形成的研究设计。例如，华盛顿州驾驶者血醇浓度高于或低于法定0.08\%界限时，所受处理截然不同。研究者据此采用断点回归设计，比较血醇略低于界限者，与恰达到或略高于界限者再次酒驾的风险，研究对酒驾处罚的作用。（Hansen B. Punishment and deterrence: evidence from drunk driving. Am Econ Rev. 2015;105(4):1581-1617.）}

\glossentry{Optimistic prior distribution}{乐观先验分布}{由于已有几项与研究问题有关的研究，认为治疗有效，而且效果通常在较窄范围内的一种先验分布。例如，研究一种类似二甲双胍的新药时，研究者使用乐观先验：它是以血糖降低30\%为中心的钟形分布，大部分分布位于降低20\%至40\%之间。}

\glossentry{Ordinal variable}{有序分类变量}{取值具有逻辑顺序的分类变量。例如，总体照护满意度，以从很差到很好的五级有序尺度测量。另见\emph{无序分类变量}。}

\glossentry{Outcome}{结局}{研究终点的总称，例如死亡或疾病发生。例如，研究放射外科治疗是否有益于单发脑转移患者时，随访的结局是死亡，或进入专业护理机构。}

\glossentry{Outcome variable}{结局变量}{对每名参与者结局的正式定义。例如，研究不同运动方式对体重及身体组成的作用，结局变量可定义为：一年后末次测量较研究开始时的体重变化，单位kg；以及同一时期的腰围变化，单位cm。}

\glossentry{Overdiagnosis}{过度诊断}{诊断出一种疾病，而如果没有作出诊断，该疾病就永远不会影响患者。例如，前列腺癌筛查的一项风险是过度诊断，而且随患者年龄增大而增加。}

\glossentry{Overfitting}{过拟合}{研究者选择多变量调整模型的变量或界点时，部分依据了样本中的偶然变异，从而使模型表现估计过于乐观的问题。例如，某模型从25个已测预测因素中选取4个，仅有20例结局事件，却报告对再次椎体骨折的预测表现极好，
\sourcepage{713}
评审者因而怀疑过拟合。}

\glossentry{Overmatching}{过度匹配}{为控制混杂而作的匹配超过必要程度，使对照与病例过于相似，降低研究者判断危险因素是否与结局有关的能力。例如，研究注射吸毒者心内膜炎危险因素时，让病例推荐同样注射毒品的熟人作对照，可能导致过度匹配，因为他们很可能有相似注射技术及毒品来源。}



\glossentry{P value}{$P$值}{假设样本所来自的人群中没有效应，即零假设正确，仅由偶然获得至少与研究所见一样大的效应（更准确地说，是至少同样大的检验统计量）的概率。例如，零假设为喝咖啡与心肌梗死无关；研究发现喝咖啡者相对不喝者的心肌梗死相对危险度为2.0，$P=0.10$，就表示如果总体确实无关联，研究中获得2.0或更大相对危险度的概率为10\%。}

\glossentry{Paired measurements}{配对测量}{以某种方式密切关联的测量，例如同一个人的两侧、同一对双胞胎的两人，或最常见的同一参与者在两个时间点，如干预前后测量。例如，研究运动项目对II型糖尿病患者血糖控制的作用时，在研究开始及运动三个月后，配对测量糖化血红蛋白水平。}

\glossentry{Pairwise matching}{成对匹配}{依据匹配变量相同或相近的取值，把一个对照与一个病例配对；较少见的是把一个暴露者与一个未暴露者配对。它用于控制该变量的混杂。例如，研究持久性有机污染物暴露是否为子痫前期危险因素时，把每名子痫前期病例，与一名未发病对照按年龄相差不超过两岁、孕前BMI相差不超过3 kg/m$^2$匹配。}

\glossentry{Partial verification bias}{部分验证偏倚}{也称验证偏倚、检查偏倚或转诊偏倚。评价检查准确度时，参与者是否接受金标准疾病验证，部分取决于待评检查本身的结果，而且分析只纳入接受金标准者，从而造成的偏倚。
\sourcepage{714}
例如，评价胸部叩诊诊断肺炎的准确度时，只纳入拍过胸片的患者，而叩诊浊音者更容易拍片，就会因部分验证偏倚而虚增叩诊灵敏度、降低特异度。}

\glossentry{Participant}{参与者}{参加研究的人。比起“受试者”，更推荐“参与者”，因为它强调入组者是在主动促进科学进步，而不只是接受实验的对象。例如，失眠新药治疗研究的参与者，是符合条件并入组的人，他们往往缺少睡眠。}

\glossentry{Participant bias}{参与者偏倚}{见\emph{回忆偏倚}。}

\glossentry{Participant identification number}{参与者识别号}{在研究数据库中识别每名参与者的字段，通常是参与者表最左边的一列。例如，一项预计纳入约250人的研究，使用从101开始的识别号，使所有人的号码都是三位数。}

\glossentry{Participant variability}{参与者变异}{参与者本身的变化，导致重复测量出现差异。例如，短暂情绪及睡眠是否充足会影响认知表现，使其测量产生参与者变异。}

\glossentry{Peer review}{同行评审}{由研究者的同行审查其方案、申请书或论文。例如，提交NIH的资助申请，由同领域科学家依据明确标准给方案评分。同样，医学期刊稿件也由科学家评审，提出改进建议，并帮助编辑决定是否发表。}

\glossentry{Per-protocol analysis}{符合方案分析}{临床试验中，仅纳入遵循方案者资料的分析方法；遵循方案通常定义为按要求服药或使用研究干预。例如，比较手术与物理治疗重度膝骨关节炎的随机试验，只分析手术组确实手术者，以及物理治疗组依从治疗方案者。另见\emph{意向治疗分析、按实际治疗分析}。}

\sourcepage{715}
\glossentry{Person-time}{人时}{研究或人群中，每名参与者处于风险中的时间之和，是计算发病密度的分母；等于处于结局风险中的人数乘以平均风险时间。例如，1000名参与者平均处于风险中2.5年，总随访人时就是2500人年，即使其中5\%只随访了一个月或更短。}

\glossentry{Phase I trial}{I期试验}{在少量人类自愿者中，逐渐增加新治疗剂量以检验安全性的早期试验，通常无盲法、无对照。例如，一种治疗绝经期潮热的新药I期试验，通常纳入少量有或无潮热的自愿者，给予递增剂量，观察血细胞计数、肝肾功能、体征、症状及其他意外不良事件。}

\glossentry{Phase II trial}{II期试验}{一项小型随机试验，最好采用盲法，检验不同剂量新治疗对副作用、生物标志物或临床结局的影响。例如，潮热新药已在I期证实安全后，II期可纳入少量有潮热的绝经后女性，随机给予两三种不同剂量或安慰剂，随访潮热频率及副作用。}

\glossentry{Phase III (pivotal) trial}{III期（关键性）试验}{规模足以检验新治疗有效性与安全性的随机试验，最好采用盲法。例如，II期已确定潮热新疗法的最佳剂量，安全性也可接受，下一步便是在大型III期中，把有潮热的绝经后女性随机分配到新治疗或安慰剂，随访潮热及不良作用。}

\glossentry{Phase IV trial}{IV期试验}{药物经FDA等监管机构批准后开展的大型研究，可以随机，也可以不随机；常用于考察比典型III期更长期的安全性。例如，潮热新药获FDA批准后，IV期可纳入潮热程度比III期入组者更轻的女性。}

\glossentry{Pilot study}{先导研究}{判断全面研究是否可行，并改进实施安排、尽可能提高效率的小型研究。例如，恢复性瑜伽预防胰岛素抵抗患者糖尿病的先导研究，
\sourcepage{716}
可能旨在说明胰岛素抵抗能够测量，改进并标准化瑜伽干预，以及证实可以招募参与者并随机分入瑜伽和对照组。}

\glossentry{Placebo control}{安慰剂对照}{随机盲法试验使用的无活性对照，理想情况下与有效药物或干预无法区分。例如，尿失禁新疗法随机安慰剂对照试验中，安慰剂应与新药外观、气味、味道及触感相同。}

\glossentry{Placebo effect}{安慰剂效应}{治疗产生了不能归于其特定性质的效果，因此很可能源于参与者对治疗的信念。例如，给慢性失眠大学生服用无活性物质后，因安慰剂效应，睡眠时间延长。}

\glossentry{Plagiarism}{剽窃}{研究者挪用他人想法、结果或文字，却不适当注明来源的一种科研不端。例如，使用另一研究者对新测量方法的描述，却不适当标明出处，就构成剽窃。}

\glossentry{Platform trial}{平台试验}{使用单一试验基础设施和主方案，同时或依次评价多种干预的适应性临床试验。例如，I-SPY 2平台试验比较一系列新药联合标准化疗，与单用标准治疗的乳腺癌疗效。}

\glossentry{Polychotomous categorical variables}{多分类变量}{具有三个或更多类别的分类变量。例如，ABO血型包括A、B、AB和O型，属于多分类变量。}

\glossentry{Population}{人群}{具有指定特征的所有人的集合。例如，可从美国全体十一年级学生中抽样，估计经常使用电子烟的比例。}

\glossentry{Population-based sample}{基于人群的样本}{能够代表整个人群的样本。例如，美国国家健康与营养检查调查（NHANES）提供美国全体人口随机样本的资料，属于基于人群的样本。}

\glossentry{Positive predictive value}{阳性预测值}{见\emph{预测值，阳性}。}

\sourcepage{717}
\glossentry{Postaward manager}{资助后管理人员}{确保研究者从资助或合同得到的经费用途适当的工作人员。例如，资助后管理人员每月编制研究支出及剩余经费报告。}

\glossentry{Post hoc hypothesis}{事后假设}{分析数据之后形成的假设。例如，研究睡眠质量与频繁跌倒风险，研究者查看了许多年龄和性别亚组中各种睡眠障碍的作用后，才提出“失眠仅增加75岁及以上男性、但不增加女性或较年轻男性跌倒风险”的假设，这就是事后假设。}

\glossentry{Posterior distribution}{后验分布}{依据研究结果，通常是临床试验结果，及先验分布，确定治疗各个疗效水平的可能性，常用可信区间或最高后验密度区间概括。例如，新血管紧张素受体阻滞剂对偏头痛发作频率作用的后验分布，以95\%可信区间表述为：相对研究开始时减少15\%至40\%。另见\emph{贝叶斯分析}。}

\glossentry{Posterior probability}{后验概率}{
\begin{enumerate}
\item 在诊断检查中，将先验概率，即根据临床特征初步估计的患病等概率，与额外信息，例如检查结果结合，所得的概率。例如，依据大流行期间病毒正在传播，以及患者发热、咳嗽，医师估计COVID-19先验概率为60\%。随后快速抗原阴性，根据检查特征，把后验概率修正为40\%。
\item 在研究结果的贝叶斯分析中，将先前治疗有效的可能性与临床试验结果结合所得的概率。例如，依据同一治疗既往研究，研究者事先认为，使中位生存期至少增加30\%的概率为70\%。完成试验并发现有益后，治疗使中位生存期至少增加30\%的后验概率升至95\%。
\end{enumerate}}

\glossentry{Power}{检验效能}{总体真实效应等于指定效应量时，在样本中正确拒绝零假设的概率。例如，假定运动使总体糖尿病女性的空腹血糖平均降低20 mg/dL。
\sourcepage{718}
研究者设效能90\%，反复从总体抽样，每次用相同测量开展同一研究，则大约每十次中有九次能正确拒绝零假设，认为运动降低空腹血糖。另见\emph{beta}。}

\glossentry{Practice-based research networks}{以实践为基础的研究网络}{社区临床医师共同研究感兴趣问题的网络。例如，初级保健腕管综合征治疗的一项网络研究发现，大多数患者经保守治疗好转；这与学术医疗中心既往文献认为多数需要手术的结果不同。}

\glossentry{Preaward manager}{资助前管理人员}{帮助研究者依照机构及监管要求，准备和提交资助申请、预算的工作人员。例如，资助前管理人员要求主要研究者，至少在截止两周前提供预算说明草稿，以便有时间检查金额是否符合资助方要求。}

\glossentry{Precision}{精密度}{在相同条件下重复测量变量，每次得到近似数值的程度。例如，自动细胞计数仪对中性粒细胞绝对数的估计，比通过显微镜观察白细胞更精密。另见\emph{准确度}。}

\glossentry{Preclinical study}{临床前研究}{在人体检验干预之前开展的研究，可以涉及细胞、组织或动物。例如，FDA要求新药先在两个不同动物物种中开展临床前研究，记录其安全性，之后才可进行人体检验。}

\glossentry{Predictive validity}{预测效度}{依据测量预测相关结局的能力，判断其是否代表所要测量的基础现象。例如，若某抑郁测量与后续自杀风险有关，其预测效度就得到加强。}

\glossentry{Predictive value, negative}{预测值，阴性}{检查阴性者没有目标疾病的概率。例如，62～91岁男性前列腺特异抗原（PSA）$\leq4.0$ ng/mL的阴性预测值约85\%。（Thompson IM, Pauler DK, Goodman PJ, Tangen CM, Lucia MS, Parnes HL, et al. Prevalence of prostate cancer among men with a prostate-specific antigen level < or =4.0 ng per milliliter. N Engl J Med. 2004;350(22):2239-2246.）
\sourcepage{719}
另见\emph{患病率、先验概率、灵敏度、特异度}。}

\glossentry{Predictive value, positive}{预测值，阳性}{检查阳性者患目标疾病的概率。例如，男性人群前列腺癌患病率为15\%时，PSA $>4.0$ ng/mL的阳性预测值约30\%。另见\emph{患病率、先验概率、灵敏度、特异度}。}

\glossentry{Predictor variable}{预测变量}{考虑两变量关联时，先发生或从生物学看更可能导致另一变量者。例如，考察肥胖是否与睡眠呼吸暂停风险增加有关，肥胖就是预测变量。按意向治疗分析的随机试验，预测变量是组别分配。}

\glossentry{Preliminary data}{初步数据}{研究开始前收集，用于判断研究成功可能性的数据。例如，为了解招募计划效果，研究者联系了符合纳入标准且不触及排除标准的人；20名中有15人对研究很热心，提示可能能够招满所需400人。}

\glossentry{Pre-post design}{干预前后设计}{见\emph{前后研究（试验）设计}。}

\glossentry{Pretest}{预试}{研究开始前，由工作人员对具体问卷、测量或程序进行评价，旨在评估并改善其功能、适当性或可行性。例如，测试录入与数据库管理系统时，让工作人员填写含缺失、越界及不合逻辑数据的表格，确认数据编辑系统能识别这些错误。}

\glossentry{Pretest probability}{验前概率}{见\emph{先验概率}。}

\glossentry{Prevalence}{患病率}{某一时点患有疾病或某状况的人所占比例，受发病率及病程影响。例如，系统性红斑狼疮（SLE）患病率，是特定时点患病者比例；如果发病增加，或一种不能治愈疾病的治疗延长患者生存，患病率就会上升。}

\glossentry{Primary hypothesis}{主要假设}{用于确定样本量及其他细节的一项研究的主要假设。例如，
\sourcepage{720}
研究的主要假设是：与外观匹配的安慰剂相比，一种针对C反应蛋白升高的新治疗，会降低吸烟者五年随访中的心血管事件风险。}

\glossentry{Primary key}{主键}{数据库表中识别每条记录，即每一行的字段，即列。每张表都应有主键。例如，Exam表的主键可命名为ExamID。另见\emph{外键}。}

\glossentry{Primary outcome}{主要结局}{反映主要问题、指导样本量计算，并决定实施工作优先次序的结局测量。例如，一项门诊COVID-19感染治疗随机试验，以住院或死亡为主要结局。}

\glossentry{Principal investigator (PI)}{主要研究者}{对设计和实施，以及结果分析和呈现承担最终责任的人。有些研究设共同主要研究者，各有具体职责。例如，伦理审查委员对方案有疑问，要求与研究的PI交谈。}

\glossentry{Prior distribution}{先验分布}{在研究结果未知前，对治疗各个疗效水平的可能性所作的估计。例如，依据另一种血管紧张素受体阻滞剂的既往研究，一种新血管紧张素受体阻滞剂对偏头痛发作频率作用的乐观先验为：减少$35\%\pm15\%$，这里为均值$\pm$标准差。另见\emph{贝叶斯分析、后验分布}。}

\glossentry{Prior probability}{先验概率}{
\begin{enumerate}
\item 检查之前，某人患病或具有另一属性的可能性，有时称验前概率。例如，因发热、咳嗽到急诊就医的成人，细菌性肺炎先验概率为12\%。另见\emph{后验概率}。
\item 研究结果未知前，估计的治疗效果可能性。例如，研究者估计，一种外周血管病新型药物洗脱支架，使置入后跛行风险至少降低25\%的先验概率约50\%。另见\emph{贝叶斯分析、先验分布}。
\end{enumerate}}

\glossentry{Probability sampling}{概率抽样}{通常利用随机数表或计算机算法，使总体每个成员都有明确入样机会的随机过程，从而为由样本推断总体提供充分依据。例如，
\sourcepage{721}
根据加利福尼亚州所有医院出院诊断，抽取慢性阻塞性肺疾病（COPD）患者的5\%概率样本，其观察结果应能反映再住院和死亡的危险因素。}

\glossentry{Professional conflicts of interest}{职业利益冲突}{专业声望、对某种观点的智识投入等非经济激励，可能使研究偏向预设结果。例如，一位整个职业生涯都想证明维生素B6治疗降低冠心病风险的研究者，若参加制定维生素B6摄入建议的小组，就可能有利益冲突。}

\glossentry{Propensity score}{倾向评分}{估计参与者具有某个指定预测变量取值的概率，最常见的是接受某种治疗的概率。通过匹配、分层或多变量分析等控制倾向评分，可以处理适应证混杂。研究者不直接调整所有可能与结局相关的因素，而先建立预测接受治疗的多变量模型，为每个人赋予预测治疗概率，即倾向评分；随后估计治疗与结局关联时，可把它作为唯一混杂指标。例如，研究阿司匹林与结肠癌关联时，用倾向评分调整与阿司匹林使用相关的因素。}

\glossentry{Proposal}{申请书}{为从资助机构获得经费而撰写，包含研究方案、预算及其他行政和支持信息的文件。例如，NIH评价申请书，以资助多种类型的研究。}

\glossentry{Prospective cohort study}{前瞻性队列研究}{对明确界定的队列，在研究开始时测量预测变量，随后随时间观察特定结局的设计。例如，护士健康研究以美国注册护士样本为队列，前瞻性研究女性常见病危险因素，结局包括心血管病、癌症及死亡。}

\glossentry{Protected health information}{受保护健康信息}{可以识别个人身份的健康信息。联邦健康隐私法规，因《健康保险可携带性与责任法》而称HIPAA法规，要求研究者在研究中保护姓名、病历号等信息的秘密性。例如，不应将它们存放在闪存盘，或通过普通电子邮件发送。}

\sourcepage{722}
\glossentry{Protective factor}{保护因素}{被认为与发生结局的可能性降低有关的属性。例如，较高的血清高密度脂蛋白（HDL）胆固醇水平，是心血管病的保护因素。另见\emph{暴露、危险因素}。}

\glossentry{Protocol}{研究方案}{一项研究的详细书面计划。例如，方案规定，只有能理解八年级水平英语的人，才符合入组条件。}

\glossentry{Psychometric adequacy}{心理测量适当性}{自报测量在人群中表现出良好心理测量特征的程度。例如，研究者确认问卷在目标人群有较高重测信度及构念效度，从而确认其心理测量适当性。另见\emph{心理测量特征}。}

\glossentry{Psychometric characteristics}{心理测量特征}{说明自报测量是否适当、相关和有用的表现特征，常包括变异性、信度、效度及常模建立。例如，研究者未在跨性别女性研究中采用女性性功能问卷，因为尚未在该目标人群评价其心理测量特征。}

\glossentry{Psychometric equivalence}{心理测量等价性}{自报工具在新或不同人群中显示相似心理测量特征的程度。例如，某抑郁问卷在大学教育受访者中有良好信度和效度，但研究者担心，对阅读能力较低者未必具有同样的心理测量表现。}

\glossentry{Publication bias}{发表偏倚}{已发表研究不能代表所有已完成研究，从而使文献失真；通常因为显示治疗有效等阳性结果，比阴性结果更常被投稿及发表。例如，一项荟萃分析只发现六项小型阳性研究及一项大型阴性研究发表，作者因而怀疑发表偏倚。}

\glossentry{Purposeful sampling (or selection)}{目的性抽样（选择）}{为了探究感兴趣的主题，以系统而非概率的方式选择参与者。例如，研究以患者为中心的医疗之家时，有目的地选择诊所管理者及一线人员访谈，考察政策如何实施及撤除。}

\sourcepage{723}
\glossentry{Quality control}{质量控制}{确保入组、测量、实验室程序及数据管理分析等研究实施工作达到最高质量的过程。例如，研究者在操作手册中为所有测量写明具体程序，并间歇观察工作人员是否遵守，以控制资料收集质量。}

\glossentry{Quality improvement}{质量改进}{医疗提供者及管理者为改善照护质量而作的努力，尤其为非研究目的实施时，可能不需要知情同意或IRB批准。例如，医院可设置自动提醒，要求定期更换中心静脉导管敷料，以降低导管感染。}

\glossentry{Quasi-experimental (quasi-randomized) design}{准实验（准随机）设计}{一种方法并不特别周密的设计，在实施干预前后测量结局，比较干预作用。例如，研究者在给小学生平板电脑前后比较阅读能力。另见\emph{前后研究}。}

\glossentry{Query}{查询}{要求关系数据库选择或操作数据的命令或指令。例如，协调员编写查询，选出未来两个月到期、却尚未安排随访的全部参与者姓名及联系方式。}

\glossentry{Questionnaire}{问卷}{通过一系列问题取得参与者信息的测量工具，可以自填，也可由工作人员施测。例如，2014年Block食物频率问卷询问127种食物的通常摄入量，以评估多种营养素及食物类别的摄入。}

\glossentry{Random error}{随机误差}{偶然变异使测量或估计偏离真实值，可通过重复测量取均值及增加样本量减少。例如，总体冠心病患者鱼油使用率确为20\%，纳入100人的研究可能恰好得到20\%，但单凭随机误差，也很可能略高或略低。}

\glossentry{Random sample}{随机样本}{列出可接触人群的单位，再随机选择其中一部分所得的样本。例如，在研究者诊所的白内障患者中随机抽样，
\sourcepage{724}
需要先列出所有白内障患者，再用计算机随机数选择。另见\emph{概率抽样、抽样框}。}

\glossentry{Randomization}{随机化}{随机试验中，把符合条件者随机分配到某一组的过程。组数及每组概率在开始随机化前确定。虽然通常以各50\%概率分成两组，也可按预定概率分入任意组数。例如，两种治疗对安慰剂，可分三组，两有效治疗组各30\%，安慰剂40\%。}

\glossentry{Randomized blinded trial}{随机盲法试验}{以预定概率随机分配合格参与者，而且研究者、参与者和其他工作人员不知道分组的设计；结局判定也不知道分组。例如，腹泻新药试验可把合格者随机分到新药或外观相同安慰剂，通常各50\%，并使研究者、参与者及工作人员均不知道所服者是哪一种。}

\glossentry{Randomized quality improvement trial}{随机质量改进试验}{在一些场所随机实施质量改进，再比较过程及健康结局，在全面推广前确认照护改善且没有意外不良后果的研究。例如，临床领导认为，由患者的医疗服务提供者签名的鼓励接种信能提高疫苗接种率，就随机给一半合格患者寄信，再测接种率是否上升。}

\glossentry{Randomized registry trial}{随机登记试验}{在持续运行的登记系统中开展的随机试验。例如，登记项目对所有接受特定器械者随访五年，收集结局和并发症。研究者可在患者接受器械时入组，随机采用新旧手术技术，再借登记系统收集结局。}

\sourcepage{725}
\glossentry{Rate}{率}{风险随时间的度量，定义为发生结局的人数除以风险人时。例如，非酒精性脂肪性肝炎（NASH）患者肝硬化发生率为每1000人年7例。另见\emph{风险率}。}

\glossentry{Rate ratio}{率比}{暴露或治疗组结局发生率，除以未暴露或未治疗组发生率。随访时间不等时，率比优于风险比，因为它考虑处于风险中的时间差异。例如，置入中心静脉导管时使用核对表，使血流感染从每1000导管日5.9例降至3.8例，率比为$3.8/5.9=0.64$。（Wichmann D, Belmar Campos CE, Ehrhardt S, et al. Efficacy of introducing a checklist to reduce central venous line associated bloodstream infections in the ICU caring for adult patients. BMC Infect Dis. 2018;18(1):267.）}

\glossentry{Recall bias}{回忆偏倚}{参与者对危险因素暴露的回忆，受到其他因素、尤其是病例或对照身份影响的一类特定偏倚。例如，肌萎缩侧索硬化病例比对照更常回忆起杀虫剂暴露，被认为可能是回忆偏倚所致。}

\glossentry{Receiver-operating characteristic (ROC) curve}{受试者工作特征曲线}{量化诊断检查准确度，并展示不同阳性阈值下灵敏度与特异度权衡的图示方法。纵轴为真阳性率，即灵敏度，横轴为相应假阳性率，即$1-\text{特异度}$，显示多个阳性界点。曲线下面积（AUROC）概括总体准确度，从无用检查的0.5到完美检查的1.0。例如，腹部超声测阑尾直径，用于诊断5～15岁儿童阑尾炎时，AUROC约0.85。（Pedram A, Asadian F, Roshan N. Diagnostic accuracy of abdominal ultrasonography in pediatric acute appendicitis. Bull Emerg Trauma. 2019;7(3):278-283.）}

\glossentry{Record}{记录}{关系数据库表的一行，最好由主键识别，其中包括特定个人、交易、结果或事件的信息。例如，Participants表可为每名参与者保留一条记录，以StudyId为主键，出生日期及性别等其他资料作为字段。}



\glossentry{Recruitment}{招募}{识别并纳入符合研究条件参与者的过程，方法随研究性质而异。例如，
\sourcepage{726}
某项研究的招募包括在专科门诊识别合格患者，以及通过传单、报纸和社交媒体广告招募。}

\glossentry{Recursive partitioning}{递归分割}{根据结局风险对人分类的多变量分析技术。与logistic回归等需要模型的技术不同，它不必假定预测与结局的关系形式，而建立一棵由一系列是／否问题组成的分类树，称分类与回归树（CART）。例如，研究者据此发现，20～65岁的急诊腹痛患者，若无食欲减退、发热及反跳痛，急性阑尾炎风险较低。另见\emph{临床预测模型、过拟合}。}

\glossentry{Referential integrity}{参照完整性}{数据库的一种性质，禁止建立所谓“孤儿记录”，即一对多关系“多”侧的记录，在“一”侧没有对应记录；也禁止删除“一”侧记录，除非“多”侧全部对应记录已经删除。例如，研究者试图从Participants表删除243号参与者，因该人在其他表尚有资料而收到错误提示。他没有感谢这一保护，反而气恼地说：“该死的参照完整性！”}

\glossentry{Registry}{登记系统}{记录某种疾病患者或接受某种程序者的数据库。可在登记过程中收集结局，或将资料与癌症登记、全国死亡索引等其他来源关联来开展研究。例如，旧金山乳腺X线摄影登记系统，收集当地三个最大检查中心所有受检女性的资料；研究者把它与当地癌症登记关联，估计乳腺摄影的准确度。}

\glossentry{Regression discontinuity design}{断点回归设计}{当某个“运行变量”决定或强烈影响是否治疗或暴露，而且越过某界点后，接受治疗的可能性明显增大或减小时，可以采用的机会性设计。例如，研究入住新生儿重症监护病房（NICU）对出生后六个月时母乳喂养的影响，可比较胎龄略高于或低于35周的新生儿；胎龄是运行变量，35周是可入住母婴病区而非NICU的界点。}

\sourcepage{727}
\glossentry{Regression to the mean}{均值回归}{极高或极低的离群测量值，在重复测量时倾向更接近总体均值。例如，因为收缩压高于第95百分位而被选入研究的儿童，即使尚未治疗，多数在第一次随访时血压也会更低，这就是均值回归。}

\glossentry{Relational database}{关系数据库}{由相互关联的表组成的数据库，为每一个参与者或属性，记录多个测量。如果某张表有一个字段，见外键，能把重复测量连回相应字段，就称其处于一对多关系的“多”侧。例如，第19章的婴儿黄疸数据库是关系数据库。}

\glossentry{Relative risk}{相对危险度}{见\emph{风险比}。}

\glossentry{Reliability}{信度}{见\emph{精密度}。}

\glossentry{Representative sample}{代表性样本}{研究入组者能够代表可接触人群的样本。例如，研究者先取得某医疗保险计划所有PSA升高男性的名单，再随机抽取5\%回顾医疗索赔资料，以获得该可接触人群的代表性样本。}

\glossentry{Reproducibility}{可重复性}{同一特征或现象未改变时，重复测量得到相同结果的程度。例如，婴儿秤每次称量标准3 kg砝码，都显示3.000 kg，就有很高可重复性。但可重复性不等于准确度：如果每次都显示3.047 kg，可重复性也同样很高。}

\glossentry{Reproducibility study}{可重复性研究}{以测量可重复性为主要问题的研究，通常比较同一人或机器多次测量的结果，即观察者内可重复性，或不同人或机器测同一对象的结果，即观察者间可重复性。例如，两名研究助理分别编码开放式问题的回答，以研究其编码可重复性。}

\glossentry{Research hypothesis}{研究假设}{研究者概括研究主要组成的陈述，包括目标或可接触人群、预测与结局变量，以及预期结果。为用于统计，假设应写成可进行显著性检验的形式，
\sourcepage{728}
通常包括零假设及备择假设。例如，研究假设为：与放射治疗相比，采用前列腺切除术治疗前列腺癌，与尿失禁风险至少增加40\%有关。}

\glossentry{Research misconduct}{研究不端}{见\emph{科研不端}。}

\glossentry{Research proposal}{研究申请书}{见\emph{申请书}。}

\glossentry{Research question}{研究问题}{研究项目试图回答的问题。好的问题应包括所关注的预测因素、结局和研究人群，通常为“在人群C中，A是否与B有关？”；临床试验则可为“在人群C中，A是否引起B？”。例如：“糖尿病成人经常使用牙线，是否降低冠状动脉事件风险？”}

\glossentry{Respect for persons}{尊重个人}{承认所有人均有权自行决定是否参加研究的基本伦理原则。它要求取得自愿知情同意；对同意能力受损者加以保护；允许随时退出；并保护隐私。例如，塔斯基吉梅毒研究未告知参与者正在接受研究，也未取得知情同意，违反了这一原则。}

\glossentry{Response rate}{应答率}{符合条件者中，回答问卷或其中某一问题的人所占比例。低应答率可能降低内部效度，使结果有偏倚。例如，高中生调查中，大麻使用问题的应答率仅20\%，就提示结果不大可能有效估计真实使用率。另见\emph{缺失数据}。}

\glossentry{Retrospective}{回顾性}{字面意思为“向后看”，用法不一致，而且常带贬义，指研究者决定开展研究或测量预测变量之前，结局已经发生的研究。例如，研究共用帽子是否为头虱危险因素的回顾性病例对照研究，可向已诊断头虱儿童家庭及未被诊断者的对照样本寄发共用帽子问卷。}

\glossentry{Retrospective cohort study}{回顾性队列研究}{队列建立、初始测量及随访均已在过去发生的队列研究。例如，为描述胸主动脉瘤自然史，
\sourcepage{729}
研究者在2022年取得2017年新诊断患者的影像记录，再用出院记录及全国死亡索引，确定哪些人在2022年前发生主动脉瘤破裂或死亡。}

\glossentry{Risk}{风险}{事件发生的概率。例如，美国每年遭雷击风险约为五十万分之一（\url{https://www.cdc.gov/disasters/lightning/victimdata.html}）。见\emph{累积发病率}。}

\glossentry{Risk difference}{风险差}{一组结局风险减去比较组风险。例如，正在使用雌激素者静脉血栓栓塞风险为$5/1000$，即0.5\%，从未使用者为$2/1000$，即0.2\%，则使用者相对未使用者的风险差为$3/1000$，即0.3个百分点。另见\emph{需伤害人数、需治疗人数}。}

\glossentry{Risk factor}{危险因素}{被认为与结局发生可能性增加有关的属性。例如，女性是多发性硬化的危险因素。另见\emph{暴露、保护因素}。}

\glossentry{Risk model}{风险模型}{假定暴露如何影响结局风险的方式。常见的是加法模型，即暴露对风险差有一致影响；或乘法模型，即暴露对风险比、率比、风险率比及比值比有一致影响。例如，研究者以乘法风险模型logistic回归，研究年龄、运动、腿部力量、饮酒及安眠药对损伤性跌倒的风险，发现饮酒使比值乘以2.7，约相当于年龄增加五岁的作用。}

\glossentry{Risk ratio (relative risk)}{风险比（相对危险度）}{一组结局风险除以比较组风险。例如，正在用雌激素者静脉血栓栓塞风险为$5/1000$（0.5\%），从未使用者为$2/1000$（0.2\%），使用者相对未使用者的相对危险度为2.5。另见\emph{风险率比、比值比}。}

\glossentry{Risk set}{风险集}{采用发病密度抽样的病例对照研究中，某病例被诊断时，尚未成为病例、仍处于结局风险中，因而可以选作对照的一组参与者。例如，研究妊娠期口服氟康唑是否影响自然流产，采用发病密度病例对照设计；
\sourcepage{730}
每发生一例流产，风险集包括妊娠已达到同一时点、尚未流产的女性。（Bérard A, Sheehy O, Zhao JP, et al. Associations between low- and high-dose oral fluconazole and pregnancy outcomes: 3 nested case-control studies. CMAJ. 2019;191(7):E179-E187.）}

\glossentry{Run-in period}{导入期}{临床试验中，让符合条件者先服安慰剂或接受有效干预的一小段时间；只有达到指定依从性、能耐受干预，或某中间结局有改善者，才能进入主要试验。例如，心律失常抑制试验中，只有导入期服有效药后室性期前收缩充分减少者，才随机分到继续药物或改安慰剂。}

\glossentry{Sample}{样本}{可接触人群中参加研究的子集。例如，哮喘新疗法研究的目标人群为所有哮喘儿童，可接触人群为今年研究者所在城镇的哮喘儿童，样本就是其中今年入组者。}

\glossentry{Sample size}{样本量}{可以指为使研究成功所估计需要的参与者人数，也可以指实际入组人数。例如，为有90\%效能发现接触暴力电子游戏使三年级男生攻击行为风险加倍，研究者估计需54人；实际样本量为58名儿童。}

\glossentry{Sampling}{抽样}{符合条件人数多于估计所需样本量时，选择参与者入组的过程。例如，研究者掷骰子，平均选入六分之一，即掷出6者。另见\emph{整群抽样、连续样本、便利样本、概率抽样、分层随机抽样、系统样本}。}

\glossentry{Sampling bias}{抽样偏倚}{见\emph{选择偏倚}。}

\glossentry{Sampling frame}{抽样框}{可接触人群中所有可能被纳入者的名单。例如，抽样框由Topeka医疗保险计划中，至少18岁、上一年被诊断发作性睡病的全部成员组成。}

\glossentry{Scale}{量表}{通过多个问题，把评分合并为量表，以测量抽象概念的常用方式。例如，测量生活质量的SF-36问36个问题，
\sourcepage{731}
产生八个功能健康与福祉有关的分量表；SF是“简式”（short form）的缩写。另见\emph{Likert量表}。}

\glossentry{Scientific misconduct}{科研不端}{有意欺骗科学界的行为总称，包括研究不端，即捏造和篡改资料、剽窃，以及挂名或隐名作者、未披露或未管理的利益冲突。例如，研究者未披露自己在所研究器械制造公司的股权，所在机构因此判定科研不端。}

\glossentry{Secondary data analysis}{二次数据分析}{用已有资料研究不同于原收集目的的问题。数据来源可包括既往研究、病历、医疗计费资料及死亡证明。例如，原来研究多危险因素干预是否减少冠心病风险的资料，被用于二次分析，考察筛查人群血清胆固醇与死亡的关联。}

\glossentry{Secondary hypothesis}{次要假设}{研究的附加假设，通常不止一个；相对于主要假设，研究对其检验的效能可能较低。例如，次要假设为：与外观匹配安慰剂相比，一种针对C反应蛋白升高的新治疗，会降低吸烟者五年随访的心血管死亡风险。}

\glossentry{Secondary research question}{次要研究问题}{主要问题之外的问题，常涉及额外预测或结局。例如，主要问题若为孕妇饮酒与低出生体重的关联，次要问题可为饮酒与妊娠期贫血的关联。}

\glossentry{Selection bias}{选择偏倚}{使选入研究的样本不能代表目标人群的系统误差。例如，在硬皮病预后因素研究中，如果参与者全都曾在学术医疗中心接受照护，就可能有抽样偏倚，因为这些患者通常病情更重，也可能比典型患者更容易获得照护。}

\glossentry{Selection criteria}{选择标准}{参与者取得入组资格所须具备的一组属性，包括纳入与排除标准。如果纳入不同组，例如病例对照或双队列，
\sourcepage{732}
标准可能不同。例如，痛风新药研究要求20～75岁、过去12个月至少一次经医师诊断的发作、血尿酸至少6 mg/dL，且无类似药物过敏史。}

\glossentry{Self-reported measure}{自报测量}{参与者直接报告所关注现象的工具，例如问卷、结构式访谈、日记或日志。例如，研究者不愿采用自报体重，因为担心有人不知道真实体重，或会低报。}

\glossentry{Sensitivity}{灵敏度}{患有目标疾病者中检查阳性的比例，即“患病者中的阳性”（positive in disease，PID）。通常病情越晚期，灵敏度越高。例如，以活检病理为参照，PSA $\geq4.0$ ng/mL检出前列腺癌的灵敏度约20\%，即每五名患者仅一名达到该浓度；检出高级别前列腺癌的灵敏度约50\%，即约一半患者达到。另见\emph{似然比、阴性预测值、阳性预测值、特异度}。}

\glossentry{Sensitivity analysis}{敏感性分析}{用不同方法，例如改变预测或结局定义、采用不同统计检验，判断主要分析结果是否经得起这些改变。例如，选择性5-羟色胺再摄取抑制剂治疗抑郁的临床试验荟萃分析，可以只纳入盲法且随访率至少90\%的试验，考察限制为高质量研究后，结果是否仍然成立。}

\glossentry{Sensitivity to change}{变化敏感性}{工具检出所测构念有意义变化的能力。例如，一份更短的焦虑新问卷变化敏感性有限：即使标准工具显示参与者焦虑严重度明显改善，新问卷评分也几乎不变。}

\glossentry{Shared effect}{共同结果}{由多个原因造成的结果。它很重要，因为以共同结果为条件，即调整它，可能在其原因之间引入虚假关联。例如，早孕与胃肠炎均可引起呕吐。
\sourcepage{733}
如果研究年轻女性呕吐原因，可能发现早孕似乎能预防胃肠炎。}

\glossentry{Simple hypothesis}{简单假设}{只有一个预测与一个结局变量的假设。例如，把维生素对胃肠道疾病影响的复杂假设，改写为“25-羟维生素D低于20 ng/mL者更容易发生憩室炎”。}

\glossentry{Simple random sample}{简单随机样本}{从抽样框中随机选择，而且每人入样概率相同的样本。例如，找出过去六个月曾到当地急诊就医的所有5～10岁儿童，为每人分配0～1随机数，再选择数值不超过0.5者。}

\glossentry{Skeptical prior distribution}{怀疑性先验分布}{认为治疗不大有效，而且效果通常在较窄范围内的先验，常因尚无人类相关研究。例如，研究新型降糖药时，采用以血糖降低0\%为中心的钟形先验，范围从降低20\%至升高20\%。}

\glossentry{Snowball sampling}{滚雪球抽样}{请参与者及利益相关者推荐其他研究对象的方式。例如，访谈一名护士后，请她推荐同病区其他工作人员参加。}

\glossentry{Specific aims}{具体目标}{申请书中对研究目标的简短陈述。例如，男性睾酮与骨密度随机试验的一项目标可以是：“检验下列假设：与随机分配接受安慰剂贴剂的男性相比，随机分配接受睾酮贴剂的男性在三年治疗期间骨丢失更少。”}

\glossentry{Specification}{限定}{在设计阶段，把某混杂因素的一个取值规定为纳入标准，以处理混杂的策略。例如，研究安抚奶嘴与婴儿猝死综合征时，只纳入配方奶喂养婴儿。若奶嘴使用者猝死风险较低，就不可能因为他们更常接受母乳。}

\sourcepage{734}
\glossentry{Specificity}{特异度}{无目标疾病者中检查阴性的比例，即“健康者中的阴性”（negative in health，NIH）。例如，以活检病理为参照，PSA $>4.0$ ng/mL检出前列腺癌的特异度约90\%，也就是90\%无癌男性PSA $\leq4.0$ ng/mL。另见\emph{似然比、阴性预测值、阳性预测值、灵敏度}。}

\glossentry{Spectrum bias}{疾病谱偏倚}{样本中的疾病谱影响灵敏度，非疾病谱影响特异度；它们若与预定应用人群不同，就会使样本中准确度不同于人群准确度。例如，一种新血清检查，以晚期食管癌与健康医学生比较时较准确，但在尚未诊断、吞咽困难的老年患者中表现很差。}

\glossentry{Split sample validation}{分样本验证}{为避免过拟合，把样本分为两组，在一组推导临床预测规则或其他模型，再在另一组验证的方法。例如，用出生住院期间可得资料，预测新生儿日后发生的黄疸；在80\%的新生儿中建立规则，再在其余20\%中验证。}

\glossentry{Standard deviation}{标准差}{连续变量变异或散布的度量，等于方差平方根。例如，400名男性队列平均年龄59岁，标准差10岁。}

\glossentry{Standard error (of the mean)}{均值标准误}{样本连续变量均值精度的估计，取决于标准差及样本量的平方根。95\%置信区间大约从均值向两个方向各延伸两个标准误。例如，400名男性平均59岁、标准差10岁，标准误为0.5岁。}

\glossentry{Standard operating procedures}{标准操作规程}{描述如何实施研究的一套或多套书面指南。例如，临床研究的方案、操作手册、统计分析计划及数据与安全监测计划，概述其标准操作规程。}



\sourcepage{735}
\glossentry{Standardization}{标准化}{为测量制定具体详细的操作指示，以尽可能提高可重复性及精密度。例如，血压测量标准化可规定如何准备参与者、使用何种大小袖带、放在何处、充气至多高及如何放气，以及哪些声音标志收缩压和舒张压。}

\glossentry{Standardized effect size}{标准化效应量}{在样本量规划中，比较两组连续变量时使用的术语，定义为预期两组均值差除以其标准差，没有单位。例如，比较80～89岁与90岁及以上者的血红蛋白，预计均值分别为13.7和13.2 g/dL，两组标准差均为1.2 g/dL，则标准化效应量为$0.5/1.2=0.42$。}

\glossentry{Statistical program}{统计程序}{用于分析和可视化数据的程序，例如Stata、SAS、SPSS或R。与电子表格不同，这些程序的列有专门的变量名、标签、类型及格式，并提供编程语言，允许交互输入命令，或把一系列命令储存为程序，日后执行。例如，研究者因为R向公众免费，而决定用它分析资料。}

\glossentry{Statistical significance}{统计学显著性}{对比较两组或更多组的研究结果作出的一项判断：若组间无差异的零假设为真，该结果不大可能出现，因此随机误差不太可能解释发现。通常把结果的$P$值，准确说相应检验统计量的$P$值，与预设阈值alpha比较。例如，随机分配接受非酒精性脂肪性肝炎新治疗的参与者，其肝病进展少于安慰剂对照组，$P=0.02$，事先alpha为0.05，研究者就可认为达到显著性，并推断治疗有效。}

\glossentry{Statistical test}{统计检验}{评价数据，为拒绝或未能拒绝某个特定假设，通常是零假设，提供定量依据的方法。连续、二分类等不同资料采用不同检验。例如，
\sourcepage{736}
研究者用卡方检验，比较认为COVID-19疫苗有效的男性与女性比例。}

\glossentry{Steering committee}{指导委员会}{为多中心研究提供整体治理的委员会，通常包括各研究地点的主要研究者、协调中心研究者及工作人员，以及资助方代表。例如，指导委员会决定拟议附属研究是否应开展。}

\glossentry{Stepped wedge trial}{阶梯楔形试验}{整群随机化的一种变体：群组或一组群组随机分配的并非干预还是对照，而是开始干预的次序。初始资料收集之后，群组按固定时间间隔，即阶梯，随机开始干预，一直持续到试验结束。例如，用这种设计在一家专业护理机构的12个病区中，研究鼓励洗手的可穿戴设备。}

\glossentry{Stratification}{分层}{分析阶段控制混杂的策略：按潜在混杂因素水平把参与者分为各层，分别分析预测与结局关联。例如，运动不足可能与骨关节炎风险增加有关，因为不运动者许多肥胖，而肥胖本身增加风险。为减少这种混杂，按研究开始时BMI分为正常、超重及肥胖层，分别分析。}

\glossentry{Stratified blocked randomization}{分层区组随机化}{旨在确保具有某种特征，通常是混杂因素，的参与者人数相等地随机进入各组的方法。先按特征分层，再在层内按预定大小区组随机。例如，既往椎体骨折既强烈预测以后骨折，也预测多种治疗反应；在预防骨折试验中，最好使有既往椎体骨折的参与者在各研究组之间人数相等，无既往椎体骨折者在各组之间的人数也相等。因此，研究者分为有和无椎体骨折两层，在每层按6～10人的区组随机。}

\sourcepage{737}
\glossentry{Stratified random sampling}{分层随机抽样}{按年龄、种族、性别等特征把潜在参与者分层，再从每层随机抽取样本。各层可按不同方式赋权。例如，研究者在加利福尼亚州胰腺癌患病率研究中，用此法过量抽取种族和族裔少数群体。}

\glossentry{Study plan}{研究计划}{说明拟议研究的基本组成，包括假设、设计、样本、主要测量，以及适用时的干预，以确定可行性、重要性、新颖性并保证合乎伦理。例如，“营养是否影响老年人肌肉功能”，可发展为随机安慰剂试验：在80岁及以上健康者中，检验每天1600 IU维生素D、持续四周对股四头肌功率的影响。}

\glossentry{Study start-up}{研究启动阶段}{首位参与者入组或首项资料收集之前的工作时期。随设计不同，可包括定预算、拟订签署合同、规定岗位、招聘培训、取得IRB批准、编写操作手册、开发测试表格及数据库，以及安排招募。例如，研究者低估了首人入组前的启动工作，使招募延迟超过三个月。}

\glossentry{Subgroup analysis}{亚组分析}{在研究参与者的某个子集中进行组间比较。例如，在选择性雌激素受体调节剂（SERM）对乳腺癌复发作用的随机试验中，分别在I、II及III期乳腺癌女性中比较SERM与安慰剂的效果。}

\glossentry{Subject}{受试者}{见\emph{参与者}。}

\glossentry{Suppression}{抑制}{混杂因素与预测变量有关，却以相反方向影响结局，从而减小表面预测与结局关联的一类混杂。例如，如果吸烟者较年轻，却未控制年龄混杂，吸烟与皮肤皱纹的关联可能被漏掉，即被“抑制”。}

\glossentry{Surrogate marker (or surrogate outcome)}{替代标志物（替代结局）}{与有意义的临床结局相关，且有时被用来代替它的测量，
\sourcepage{738}
因为治疗引起的指标变化出现更快，或发生在更多参与者中。替代指标应中介治疗对结局的大部分作用。例如，HIV感染者CD4淋巴细胞数增加，预示机会性感染风险降低，常用作抗逆转录病毒药疗效的替代标志物。但也有潜在替代指标改善、临床结局却未改善的例子。托塞曲普虽然有利地改变血清低密度及高密度脂蛋白胆固醇，而两者常被视为动脉粥样硬化进展的替代标志物，却增加了死亡及心血管发病。（Barter PJ, Caulfield M, Eriksson M, et al. Effects of torcetrapib in patients at high risk for coronary events. N Engl J Med. 2007;357:2109-2122.）另见\emph{中间标志物}。}

\glossentry{Survey}{调查}{在特定人群中进行、通常使用问卷的横断面研究。例如，全国酒精及相关疾病流行病学调查纳入美国成人代表性样本，询问当前及既往饮酒、酒精使用障碍，以及酒精治疗服务利用情况。}

\glossentry{Survival analysis}{生存分析}{比较研究各组至某结局时间的统计技术，结局不必是生存。例如，冠状动脉搭桥与经皮冠状动脉成形术预防心肌梗死及死亡的随机试验，可用生存分析比较两组从开始治疗至任一结局的时间。}

\glossentry{Systematic error}{系统误差}{见\emph{偏倚}。}

\glossentry{Systematic review}{系统综述}{系统识别某问题全部研究，采用明确纳入标准及标准化资料提取方法的医学文献综述，也可包括结果的荟萃分析。例如，研究者系统综述了所有检验补锌是否降低感冒风险的研究。}

\glossentry{Systematic sample}{系统样本}{先规定抽样框，再按预先规定、但非随机的过程选择子集所得的样本。例如，研究者找出过去六个月曾到当地急诊就诊的全部5～10岁儿童，再每隔一人选一人。
\sourcepage{739}
与简单随机样本相比，这种方法没有优点，反而有几个缺点。}

\glossentry{t test (or Student’s t test)}{$t$检验（Student $t$检验）}{判断两组连续变量均值是否显著不同的统计检验。例如，两组接受不同抗抑郁药，可比较治疗后平均抑郁评分，即非配对双样本$t$检验；也可比较从开始至治疗后的平均变化，即本书所称配对双样本$t$检验。另见\emph{单样本$t$检验、双样本$t$检验}。}

\glossentry{Target population}{目标人群}{由临床及人口学特征界定、研究者希望向其推广结果的较大人群。例如，在自己医院研究儿童哮喘新疗法，其目标可能是全世界哮喘儿童。}

\glossentry{Test-result-based sampling}{按检验结果抽样}{诊断准确度研究中，分别抽取阳性、阴性等不同结果者，使各结果的比例由研究者设定，而非自然形成的方法，常用于节省后续检查费用。例如，评价尿试纸诊断尿路感染的灵敏度及特异度，对所有阳性者作尿培养金标准，但阴性者只随机取20\%培养。}

\glossentry{Thematic coding}{主题编码}{用标签或编码标记定性资料片段，说明其与研究概念的关系；这些编码通常汇入编码手册。例如，分析者通过主题编码，识别受访者谈论健康素养的段落。}

\glossentry{Thematic (or data) saturation}{主题（资料）饱和}{继续收集和分析资料，已不再获得新的概念或理论认识的状态。例如，最后三次访谈编码未产生新主题，研究者因而认为已达到主题饱和。}

\glossentry{Time series design}{时间序列设计}{在每名参与者或整个社区接受干预前后测量的一种组内设计。每个人都是自身对照，因此消除大部分混杂，但仍易受成熟效应、均值回归及长期趋势影响。例如，
\sourcepage{740}
在一组糖尿病患者运动项目开始前后测空腹血糖，判断运动是否降低血糖。另见\emph{组内设计}。}

\glossentry{Translational research}{转化研究}{旨在把科学发现转化为健康改善的研究。可在患者临床研究中检验实验室基础发现，常称“从实验台到病床”；也可把临床发现推广到人群健康，常称“从病床到人群”。例如，研究导致小鼠先天性耳聋的遗传缺陷在人是否有类似影响，属于前者；全州使用测皮层听觉反应的新生儿筛查，研究发现听力障碍是否改善以后学校表现，属于后者。}

\glossentry{Treatment}{治疗}{见\emph{干预}。}

\glossentry{Triangulation}{三角互证}{比较访谈、观察、定量资料等多种资料，以丰富发现的细节，提高准确性或可信度。例如，完成访谈和直接观察后，研究者用三角互证，更全面理解手术室不同成员如何响应术前暂停核对流程。}

\glossentry{Truth-telling}{如实告知}{研究伦理基本原则，要求说真话，不隐瞒重要相关信息，不夸大利益，也不低估风险。例如，向癌症新疗法潜在参与者解释同意时，只强调可能利益、淡化风险讨论，就违反如实告知。}

\glossentry{Two-sample t test}{双样本$t$检验}{比较一个样本与另一个样本连续变量均值的统计检验。例如，研究者发现，橄榄油补充剂组研究期间HDL胆固醇平均增加10 mg/dL，安慰剂组增加2 mg/dL，双样本$t$检验$P=0.14$。另见\emph{单样本$t$检验}。}

\glossentry{Two-sided hypothesis}{双侧假设}{研究者同时考虑两个可能方向I类错误的备择假设，例如治疗可能使结局风险升高，也可能降低。例如，
\sourcepage{741}
检验抗焦虑药与痴呆风险增加或减少是否有关。另见\emph{单侧假设}。}

\glossentry{Type I error}{I类错误}{研究结果有统计学显著性，即$P\leq\alpha$，却错误拒绝了真实零假设。例如，膳食胡萝卜素与结肠癌研究设alpha 0.05，结果以$P<0.05$认为它降低风险，但实际无关联，就发生I类错误。另见\emph{假阳性结果}。}

\glossentry{Type II error}{II类错误}{研究中未拒绝实际上为假的零假设，即$P>\alpha$。例如，胡萝卜素实际上降低结肠癌风险，研究却以$P>0.05$未能拒绝“无作用”假设，就是II类错误。另见\emph{假阴性结果}。}

\glossentry{Undue influence}{不当影响}{使参与者违背意愿或更好的判断而参加研究的压力情形，例如报酬过高，或招募囚犯、研究者自己的学生。例如，以提前假释作为参加研究的条件，是不当影响。}

\glossentry{Uninformative prior distribution}{无信息先验分布}{不对治疗效果可能大小表达特定信念，认为给定范围内各效果同样可能的先验，通常因为没有相关研究。例如，研究新降压药对血清肌酐的作用，使用矩形先验，从升高40\%延伸至降低40\%。}

\glossentry{Validity}{效度}{测量代表所关注现象的程度。例如，睡眠紊乱问卷评分，能够在多大程度上反映睡眠实验室记录等方式所确定的睡眠紊乱，就具有多大效度。}

\glossentry{Variability}{变异性}{测量散布的大小，通常用标准差表示。例如，某饮食导致的体重变化，从明显增重到明显减重都有，就具有较大变异性。另见\emph{方差}。}

\glossentry{Variable}{变量}{能够有不同取值的测量。例如，注射毒品使用可取从不、仅过去、目前等不同值，因而是变量。另见\emph{分类变量、混杂变量、连续变量、计数变量、二分类变量、离散变量、
\sourcepage{742}
无序分类变量、有序分类变量、结局变量、预测变量}。}

\glossentry{Variance}{方差}{总体或样本连续数值散布的度量。先求均值，再把每个数值与均值差的平方相加，除以数值个数；样本方差则除以个数减一。方差不十分直观，但其平方根就是标准差。例如，一支男子职业篮球队的身高方差为92 cm$^2$，相当于标准差9.6 cm。}

\glossentry{Verification bias}{验证偏倚}{见\emph{部分验证偏倚}。}

\glossentry{Visual analog scale}{视觉模拟量表}{通常用一条线表示从一个极端至另一极端的连续答案范围。线一般长10 cm，得分是距最低端的厘米数。例如，疼痛严重度量表一端写“无痛”，另一端写“难以忍受”，参与者在最符合自己疼痛的位置标“X”。}

\glossentry{Vulnerable research participants}{易受伤害的研究参与者}{研究中更容易被不合伦理地利用的人，包括儿童、囚犯，以及难以理解利弊或易受不当影响的人。例如，重度痴呆患者属于易受伤害参与者。}

\glossentry{Wait-list design}{等待名单设计}{把参与者随机分为试验开始即接受干预者，及经过规定时间后才接受干预的等待名单对照组。例如，用此设计研究盆底肌训练对尿失禁频率的作用。}

\glossentry{Washout period}{洗脱期}{交叉研究中两种治疗之间，或治疗与对照之间预留的一段时间，用于让首个治疗的作用消退，使结局测量可恢复到研究开始的水平。例如，比较利尿药与安慰剂治疗高血压的交叉试验，可规定一个月洗脱。}

\glossentry{Within-group design}{组内设计}{在同一组参与者中，通常于两个或更多时间点比较测量的设计。每个人作为自身对照，消除了不随时间改变因素的混杂；但
\sourcepage{743}
仍易受学习效应、均值回归及长期趋势影响。例如，对肺结节病患者，在运动项目前后测肺活量，判断是否改善。另见\emph{组间设计、单样本$t$检验、时间序列设计}。}

\glossentry{Z test}{$Z$检验}{比较比例是否具有统计学显著差异的检验。卡方始终双侧，而$Z$检验也可用于单侧。例如，单侧$Z$检验可判断囚犯纹身比例是否显著高于非监禁者；双侧$Z$或卡方则可判断是否不同，即可能更低或更高。}
